\section{第二段创业：Palantir 想象、并购与痛苦急转}

上一章讲秒针，本章进入第二段创业。明略数据被广密注意到，原因之一是它像中国版 Palantir：用数据整合、分析平台和工程交付服务政府/企业客户。但中国 ToB 市场和美国不同，ToG、ToB、产品化、定制交付、现金流和组织能力都会改变路线。

\begin{figure}[H]
\centering
\includegraphics[width=0.94\textwidth]{palantir-to-b-logic.png}
\caption{Palantir 与中国 ToB：不是照搬 ToG，而是寻找数据、场景和交付闭环。自制概念图，依据 00:56:08--01:15:00 对谈内容整理。}
\end{figure}

\begin{knowledgebox}{读图：Palantir 是参照物，不是答案}
Palantir 的启发在于数据整合、分析平台和强交付能力；但明略的现实路径必须处理中国企业客户、行业数据、ToB 付款、定制化和组织成本。照搬商业模式往往不够，必须重新找本地闭环。
\end{knowledgebox}

\subsection{并购 Red/小红：不是只买 application，而是买数据入口}

前面讲 Palantir 只是参照物，本节看明略怎样在中国企业微信生态里寻找自己的数据入口。访谈里很重要的一段是收购小红团队。吴明辉看中的不是单纯 application，而是微伴等产品在企业微信生态里的 install base，以及由销售、客服和客户对话产生的大量 conversational data。这和他在北大读工程博士时研究 Conversational Intelligence 的方向正好重合。

\begin{figure}[H]
\centering
\includegraphics[width=0.96\textwidth]{conversational-data-acquisition.png}
\caption{Conversational Data 线索：并购不是只买应用，而是买数据入口和 AI 场景。自制概念图，依据 00:04:36--00:07:40 对谈内容整理。}
\end{figure}

\begin{knowledgebox}{读图：为什么对话数据有价值}
企业销售和客服对话不是普通聊天记录，而是客户需求、异议处理、交易推进、产品解释和服务质量的真实轨迹。它既是训练数据，也是工作流入口。
\end{knowledgebox}

\begin{importantbox}{并购的 AI 逻辑}
如果只看收入，收购一个 application 可能只是扩产品线；如果从 AI 角度看，收购一个高频工作流入口，等于拿到一条持续产生私有数据的管道。后者才是长期价值。
\end{importantbox}

\subsection{2020--2022：战场开太宽与现金流压力}

第二段创业后半程是痛苦急转。公司同时开多个战场，现金流压力变大，资本市场环境变化，创始人开始从理想主义 founder 变成会经营的 founder。这里的教育意义很强：ToB 公司如果产品线太宽、交付太重、组织太散，就会被现金流和管理复杂度惩罚。

\begin{figure}[H]
\centering
\includegraphics[width=0.96\textwidth]{two-b-company-turnaround.png}
\caption{ToB 公司痛苦急转：战场过宽、现金流压力和经营能力迫使公司重新聚焦。自制概念图，依据 01:15:00--01:45:01 对谈内容整理。}
\end{figure}

\begin{knowledgebox}{读图：急转不是失败，而是经营能力升级}
战场过宽带来成本压力，痛苦年份迫使公司重新聚焦，聚焦以后才可能靠产品和服务造血。对 ToB 公司而言，活下来本身就是能力的一部分。
\end{knowledgebox}

\subsection{资本敬畏：估值越高未必越好}

经历过战场收缩之后，本节回到资本结构。吴明辉谈到，很多创业者经历过这一轮后才真正敬畏资本。估值高、融资多，看起来减少稀释，但如果业务规模和上市路径跟不上，就会遇到回购、投资人压力和创始人困境。ToB 软件公司尤其如此：收入不到一定规模，资本市场出口会非常窄。

\begin{warningbox}{高估值不是护城河}
估值是未来现金流和资本市场预期的压缩，不是企业真实能力。融资越多，组织越容易扩张过快；估值越高，未来回购和退出压力越大。ToB 公司如果没有造血能力，高估值会变成债。
\end{warningbox}

\begin{knowledgebox}{公司阶段对照：同一家公司里的三种能力}
\begin{tabularx}{\textwidth}{p{0.20\textwidth}p{0.30\textwidth}X}
\toprule
阶段 & 主要问题 & 形成的能力 \\
\midrule
秒针阶段 & 广告监测、中立测量、Web2 数据 & 可信数据和交易标准意识。 \\
明略阶段 & ToB/ToG 数据分析、Palantir 想象、并购 & 行业数据、交付和组织整合能力。 \\
AI 阶段 & 私有数据、Agentic Model、可信执行 & 把数据和 workflow 转成智能体交付。 \\
\bottomrule
\end{tabularx}
\end{knowledgebox}

\subsection{本章小结}

第二段创业讲的是从技术公司到经营公司的转变。Palantir 想象提供方向，并购带来数据入口，但现金流和战场过宽迫使公司学会聚焦。

